\section{Introduction \& scope}\label{sec:intro}

QtRocket is an open-source model-rocket simulator: a C++23 simulation core shared by three
front ends --- a Qt6 Widgets GUI (\code{qtrocket}), a Qt-free headless REPL
(read--eval--print loop, \code{qtrocket-cli}), and a standalone OpenGL design viewer
(\code{qtrocket-visualizer}). All three executables link one static library,
\code{qtrocket\_core}, which contains the application controller and pulls in the
\code{model}, \code{sim}, and \code{utils} layer libraries beneath it
\src{CMakeLists.txt}{224}. The core describes a rocket as a composite tree of parts,
resolves each part's absolute position from stored placement intent, computes time-aware
mass, center of gravity (CG), and inertia while a motor burns, and integrates the resulting
three-degree-of-freedom (3-DOF) point-mass flight with a fixed-step fourth-order
Runge--Kutta (RK4) or adaptive Runge--Kutta--Fehlberg (RKF45) solver until a typed
termination condition ends it.

This document is the project's design reference: a rigorous exposition of how that software
works --- its layered architecture, its data model, its algorithms, and the physics
underneath them --- written so that an engineer new to the project, with or without a
rocketry background, can navigate the source and extend it without first reverse-engineering
the design from the code. Each chapter therefore explains not only what a subsystem does but
what it stores, what it derives, why the boundary sits where it does, and what the code
deliberately does not do yet. Physics is developed before it is used: a reader who has never
heard of a center of pressure will meet the concept in \cref{sec:primer} and in background
callouts before any chapter leans on it.

\subsection{Three documents, one series}\label{sec:intro-series}

This whitepaper is the third volume of a series with a deliberate division of labor. The
Architecture Review (\srcf{docs/ArchitectureReviewLatex/}) is a capabilities-and-gaps review
performed at the same revision this document describes: it measures what the program can
actually do, adversarially verifies every claim against the source, and ranks what is
missing against OpenRocket as a baseline. The Part Placement whitepaper
(\srcf{docs/PartPlacementDesignLatex/}) is the design specification for the placement
feature that now ships: the argument, written before implementation, for storing placement
as physical intent and deriving all absolute geometry through a single resolver. Those two
documents assess and specify; this one teaches. Where a companion's verified finding is
load-bearing here --- a located defect, a dormant-status verdict --- this document states it
in place with the same citation rather than re-deriving it, and it does not re-litigate the
review's rankings or the specification's rejected alternatives except where the rationale is
part of the design being taught.

\subsection{Ground truth and citations}\label{sec:intro-groundtruth}

Every statement about the code is pinned to one revision: commit \code{254cd41} (``Major
comment cleanup''), the tip of the \code{development} branch at the time of writing,
reviewed against a clean working tree on 2026-07-07. That commit includes the completed Part
Placement feature; work in flight on other branches is mentioned only where a chapter says
so explicitly. Code citations take the form \src{model/parts/Part.h}{130} and are placed
directly after the claim they ground. Line numbers are anchors into that revision, not
absolutes --- they will drift as the tree moves, but they resolve exactly at the pinned
commit.

\begin{designrule}[Ground truth policy]
Three sources are authoritative for this document: the source tree at pinned commit
\code{254cd41}, the Architecture Review (\srcf{docs/ArchitectureReviewLatex/}), and the Part
Placement design whitepaper (\srcf{docs/PartPlacementDesignLatex/}). Every other artifact
under \srcf{docs/} --- the architecture audit, the phase specifications, the implementation
plans --- is historical: the review verified that those files have drifted from the shipped
code, and this document never cites them as fact. Where any document, including this one,
disagrees with the tree, the tree wins. General physics background carries no citation and
is framed explicitly as theory (\cref{sec:intro-reading}); it is never attributed to the
code.
\end{designrule}

\subsection{The state of the project, honestly}\label{sec:intro-state}

QtRocket at this commit is a finished, well-tested 3-DOF point-mass simulator surrounded by
machinery staged for six degrees of freedom (6-DOF). What flies is small: thrust pinned
along the world vertical, ``assumed through the CM'' \src{model/RocketModel.cpp}{70},
pluggable gravity, and a single constant-$C_d$ quadratic drag term; there is no angle of
attack, no restoring moment, no recovery deployment, and every flight ends ballistically,
descending below the launch site \src{model/RocketModel.cpp}{62}. Around that core sits what
the companion review calls ``a ring of machinery that is built, unit-tested, and consumed by
nothing'': the per-part Barrowman center-of-pressure composition
\src{model/parts/Part.cpp}{231}, a torque interface that returns zero and has no caller
\src{model/RocketModel.cpp}{94}, quaternion orientation slots in the state vector
\src{sim/StateData.h}{46}, a commented-out orientation integrator \src{sim/Propagator.h}{103},
and a wind model that returns the zero vector \src{sim/WindModel.cpp}{16}. The staging is
deliberate, and the tree says so itself: ``nothing in production consumes this in P2 --
these tests pin the seam for P5'' \src{model/tests/AeroTests.cpp}{18}. This document teaches
the dormant machinery alongside the live code, at the same depth, because the dormant seams
are the shape of the next tier of the simulator and a new engineer will extend them, not
replace them.

\begin{keyinsight}[How this document marks unbuilt work]
Three status words recur, with the companion review's exact senses. \emph{Dormant} means
built and tested but with zero production consumers --- unreachable from any shipping code
path. \emph{Partial} means reachable but with named gaps. \emph{Absent} means no trace in
the tree. A \emph{seam} is an interface or capability built and tested ahead of its
consumer; QtRocket's development order is seams first, consumers later. Each chapter marks
its dormant, partial, and absent areas plainly, in place, where the surrounding machinery is
taught; \cref{sec:incomplete} consolidates all of them into a single inventory. Expositional
chapters do not editorialize about what should be built next --- that is the review's job.
\end{keyinsight}

\subsection{How to read this document}\label{sec:intro-reading}

Four conventions hold everywhere, and the code enforces or documents each of them.

\begin{itemize}
\item \textbf{SI units throughout}: meters, kilograms, seconds, newtons. The constant
  gravity model, for example, returns an acceleration of $(0, 0, -9.80665)\unit{m/s^2}$
  \src{sim/ConstantGravityModel.h}{18}.
\item \textbf{World frame: $z$ up.} Altitude is the $z$ component of position, the launch
  site is the origin, and the nominal end of flight is descending below $z = 0$
  \src{model/RocketModel.cpp}{62}.
\item \textbf{Part frame: \zfwd.} Every part is modeled in its own frame with $+z$ toward
  the nose, its fore plane at the local origin, and its body occupying
  $z \in [-L, 0]$ \src{model/parts/Part.h}{130}. When the placement resolver plants the root
  part's fore plane at the world origin, the nose tip becomes the datum for the whole
  assembly; composite CG and center of pressure (CP) are both reported against it
  (\cref{sec:placement,sec:aero}).
\item \textbf{Quaternions are stored $(x, y, z, w)$} --- vector part first, scalar last.
  \code{Quaternion} is \code{Eigen::Quaterniond} \src{utils/math/MathTypes.h}{16}, and the
  state vector's orientation members carry the ``(vector, scalar)'' annotation
  \src{sim/StateData.h}{46}. Eigen's value constructor takes $(w, x, y, z)$;
  \cref{sec:incomplete} flags that hazard where the dormant slots are inventoried.
\end{itemize}

Inline identifiers are set in colored teletype (\code{getCompositeMass}); verbatim code
fragments use real C++ syntax (\cppinline|getMass(t)|). Listings quote the source verbatim
and cite it in the caption. Acronyms are expanded at first use in each chapter. Beyond the
running text, five callout boxes carry distinct semantics:

\begin{description}
\item[Design rule (blue).] A rule the code enforces, or an invariant the design commits to
  --- quoting the in-tree comment where it binds whenever one exists.
\item[Key correctness detail (teal).] A subtlety that is easy to miss and expensive to
  rediscover: why a cache key is sound, why a clamp must exist, how a guard prevents a
  failure that once occurred.
\item[Hard, located failure (red).] A verified defect present at the pinned commit, stated
  with its mechanism and its file:line location. These are inherited from the companion
  review's adversarial verification and restated here where the affected machinery is
  taught.
\item[Rejected alternative (amber).] A design alternative that was considered and turned
  down, kept because the argument that killed it is part of the design.
\item[Physical background (dark blue).] Pure physics: theory a newcomer needs, developed
  from first principles. The separation is deliberate --- anything inside a physics box is
  background and carries no code citation; anything outside one that describes behavior is a
  statement about the code and does.
\end{description}

\subsection{Document roadmap}\label{sec:intro-roadmap}

\Cref{sec:primer} opens with physics, not code: a model-rocket flight from first principles
--- thrust, drag, stability, and the vocabulary (center of pressure, static margin, apogee,
recovery) the rest of the document uses freely. Readers who already know rocketry can skip
it. \Cref{sec:overview} maps the system: the three executables, the layer libraries and the
link graph that enforces their ordering, the application controller, and the
\code{Propagatable} bridge between model and simulation. Three chapters then build the
rocket: \cref{sec:parts} covers the composite part tree and its time-aware mass, CG, and
inertia machinery; \cref{sec:placement} the placement model --- placement intent stored as
fractional stations and seat kinds, the single resolver that turns intent into absolute
geometry, and the fail-closed overlap diagnostics gating it; \cref{sec:motors} the motor
subsystem, from thrust-curve interpolation and the reconstructed mass curve to the database
and its three ingest paths. Three more chapters fly it: \cref{sec:simcore} the propagator
loop, the two solvers, and the typed termination taxonomy; \cref{sec:environment} the
pluggable atmosphere, gravity, and geoid models; \cref{sec:aero} the aerodynamics --- the
small consumed force path and the dormant Barrowman layer beside it. \Cref{sec:interfaces}
covers the three user-facing surfaces, from the GUI's launch panel to the 33-command
scriptable REPL \src{cli/Repl.cpp}{245}; \cref{sec:persistence} the design and
motor-database file formats and their round-trip guarantees; \cref{sec:testing} the test
suites, coverage tooling, and continuous integration. \Cref{sec:incomplete} closes the body
with the consolidated inventory of everything dormant, partial, or absent.

Four appendices place QtRocket next to OpenRocket, the de facto standard open-source
model-rocket simulator. The intent is calibration, not conformance: QtRocket is an
independent project with similar goals, OpenRocket is the baseline because it defines what a
mature tool in this space offers, and the comparison measures distance from that surface ---
never compliance with it as a specification. \Cref{app:or-context} establishes the baseline
and that framing; \cref{app:or-features} compares feature surface;
\cref{app:or-methods} simulation methods; \cref{app:or-architecture} architecture.
OpenRocket-side claims are researched from that project's published documentation and
sources rather than verified file-by-line the way QtRocket claims are, and the appendices
say so. \Cref{app:glossary} collects the series' terms of art, and \cref{app:repro} records
how this document itself builds and the ground rules under which it was written.
